% Robot controller and action space of the pushing policy: what we use, in simulation and on the real FR3, why,
% and the measurements behind the chosen values. Written for the supervisors; kept short on purpose.
% Included by main.tex (% Robot controller and action space of the pushing policy: what we use, in simulation and on the real FR3, why,
% and the measurements behind the chosen values. Written for the supervisors; kept short on purpose.
% Included by main.tex (% Robot controller and action space of the pushing policy: what we use, in simulation and on the real FR3, why,
% and the measurements behind the chosen values. Written for the supervisors; kept short on purpose.
% Included by main.tex (% Robot controller and action space of the pushing policy: what we use, in simulation and on the real FR3, why,
% and the measurements behind the chosen values. Written for the supervisors; kept short on purpose.
% Included by main.tex (\input{open_questions/impedance_action_study}); uses only packages loaded there
% (amsmath, booktabs). Numbers: scripts/sweep_impedance.py, scripts/sweep_action_poses.py,
% scripts/check_push_env.py and docs/notes/2026-10-06.md to 2026-10-08.md in the project repo. Current values and
% their derivation: action_limits_problem.tex (Section sec:al).

% ======================================= BEGIN CONTENT =======================================
\section{Open question: robot controller and action space}
\label{sec:action-controller}

\paragraph{The problem.} The policy does not move the robot's joints directly. At 31\,Hz it outputs a small motion
command for the fork (``move a little this way, turn a little''), and a low-level controller turns these commands
into motor torques at 1\,kHz. The policy learns how the fork reacts to its commands, so two things must hold for it
to work on the real robot: (1) the controller in simulation must be the one that runs on the real FR3, with the same
values, otherwise the fork reacts differently after the transfer; (2) the fork should follow the commanded steps with
a small lag. This is not strictly required, since the policy could learn a lag, but with a small lag the motion is
set by the controller, which is identical in simulation and on the robot, rather than by the arm's dynamics, which
are not; action spaces that keep the tracking error small transfer better~\cite{aljalbout2024action}. The price is
agility (Section~\ref{sec:al}). This section proposes the controller, the form of the commands, and how large a step
may be.

\paragraph{Summary.} The policy commands small steps of the tool's position and orientation. They are executed by
Franka's Cartesian impedance controller, with one change in the damping, and the \emph{same} controller with the
\emph{same} values runs in simulation and on the real FR3 (Section~\ref{sec:controller-what}). Successful sim-to-real
work on Franka robots does the same~\cite{tang2023industreal,frankatwin}. In simulation the tool follows the
commands within 2\,mm (1.68\,mm measured) over the workspace at up to 10\,cm/s and 45$^\circ$/s, with the speed-up
limited to 0.08\,m/s$^2$ and 20$^\circ$/s$^2$ (Section~\ref{sec:al}; the first design, Section~\ref{sec:controller-results},
was limited to 5.5\,cm/s and 7.8$^\circ$/s). We ask the supervisors to
confirm the controller and to tell us about the lab's FR3 setup (end of this section).

\subsection{What we use}
\label{sec:controller-what}

\paragraph{Action.} Every policy step (31.25\,Hz) the policy outputs a displacement of the tool tip (at most
3.2\,mm, i.e.\ 10\,cm/s) and a rotation of the fork (rotation vector, at most 1.44$^\circ$, i.e.\ 45$^\circ$/s), both in
the robot's base frame. The step may change by at most 0.08\,mm (0.02$^\circ$) from one policy step to the next, which
limits the acceleration to 0.08\,m/s$^2$ (20$^\circ$/s$^2$) and with it the lag (Section~\ref{sec:al}). The steps are
applied to the \emph{previous target} pose, not to the measured tool pose (IndustReal's policy-level action
integrator~\cite{tang2023industreal}), so commanded steps are not lost while the tool moves freely. Within the policy
step the target moves and turns at constant rate, and its velocities $\dot x_d$, $\omega_d$ are passed to the
controller. At every control step the target is kept within 4\,mm and 3$^\circ$ of the measured tool pose, which bounds
the spring force to 4\,N and the moment to 5.2\,N\,m when the fork is blocked; while held, the target moves with the
tool and its velocity is the tool's. The fork starts horizontal; its orientation is then free.

\paragraph{Control law.} Every control step the joint torques are
\begin{equation}
  \tau = J^\top \begin{bmatrix} K_p\,(x_d - x) - D_p\,(\dot x - \dot x_d) \\
                                 -K_o\,e_o - D_o\,(\omega - \omega_d) \end{bmatrix}
       + \left(I - J^\top J^{\top+}\right)\left(k_\text{ns}(q_\text{ns} - q) - 2\sqrt{k_\text{ns}}\,\dot q\right)
       + \tau_\text{Coriolis},
  \label{eq:impedance}
\end{equation}
a spring and damper that pull the tool tip $x$ to its target $x_d$ and its orientation to the target orientation
($e_o$: orientation error), plus a weak spring that keeps the elbow of the redundant arm near its rest posture. This is the law of
Franka's example controllers~\cite{libfranka,frankaros}, with one change: Franka sets the damping to
$D = 2\sqrt{K}$, which is critically damped only for a tool of 1\,kg. The arm's inertia felt at the tool,
$\Lambda = (J M^{-1} J^\top)^{-1}$, is several kilograms and depends on pose and direction. We therefore use
$D = \Lambda^{1/2} K^{1/2} + K^{1/2} \Lambda^{1/2}$~\cite{albu2003cartesian}, i.e.\ $2\sqrt{K}\,\Lambda^{1/2}$,
critically damped everywhere. Values: $K_p = 1000$\,N/m, $K_o = 100$\,N\,m/rad (stiffer than Franka's examples,
which use up to 400\,N/m and 30\,N\,m/rad, to keep the lag small; Section~\ref{sec:al}), $k_\text{ns} = 0.5$\,N\,m/rad,
torque change at most 1\,N\,m per ms (Franka's limit).

\paragraph{Simulation and real robot.}
\begin{center}
\begin{tabular}{p{0.2\textwidth}p{0.36\textwidth}p{0.36\textwidth}}
  \toprule
  & Simulation (Isaac Lab, PhysX) & Real FR3 \\
  \midrule
  Law~\eqref{eq:impedance} & Python function, every 2\,ms, joint torques to the simulator
    & the same equations in C++ in a libfranka torque loop at 1\,kHz \\
  $J$, $M$, Coriolis & from the simulator (Coriolis left out; small below 0.1\,m/s)
    & from Franka's robot model in libfranka \\
  Gravity & switched off for the robot & compensated by the robot (fork mass entered as load) \\
  Motor inertia & added to every joint (gear ratio$^2\times$motor inertia, from Franka's description files)
    & physical; whether libfranka's $M$ includes it is to be checked \\
  Target clamp, feedforward & every 2\,ms, function \texttt{substep\_command} & the same function in C++ in the 1\,kHz loop \\
  Speed cap, acceleration limit & policy side, every 32\,ms & the same Python code, every 32\,ms \\
  Policy & same process, every 32\,ms & Python process, sends the step to the C++ loop every 32\,ms \\
  \bottomrule
\end{tabular}
\end{center}
To make sure both sides compute the same torques, the C++ version will be checked against the Python version on
recorded robot states. Remaining differences (friction, exact inertia, delays) are handled by measuring the real
robot's response to the same commanded paths and by randomizing gains and timing during
training~\cite{peng2018sim,frankatwin}.

\subsection{Why}

The literature that transferred contact-rich tasks to a real Franka runs the same task-space impedance
controller with the same gains in simulation and on the robot~\cite{tang2023industreal,frankatwin,kaspar2020sim2real},
and identifies or randomizes what still differs~\cite{peng2018sim}. The choice of action space measurably affects
transfer~\cite{aljalbout2024action}; for an overview of the reality gap see~\cite{aljalbout2026realitygap}.
The placeholder we started from (inverse kinematics and stiff joint position control with Isaac Lab's default gains)
lagged 0.2\,s behind its targets and has no counterpart on the real robot.

\subsection{Results in simulation}
\label{sec:controller-results}

\emph{These results led to the first design (speed cap only, $K_p = 400$\,N/m, $K_o = 30$\,N\,m/rad, at most
5.5\,cm/s). They are kept as history; the current values and measurements are in Section~\ref{sec:al}.}

Criterion for ``the step is followed'': while moving, the tool trails the commanded path by at most 2\,mm, and it
overshoots the end point by at most 5\,mm (fork tilt at most 1$^\circ$); while turning, the angle lags and overshoots
by at most 2$^\circ$ and the tool tip moves by at most 2\,mm. Every step is eventually executed: after a
hold the tool is within 0.5\,mm of the commanded end point in all tests.

\begin{table}[h]
  \centering
  \caption{Choices tested (fork in free space). Left: 10\,cm commanded, 5\,cm/s, $K_p = 150$\,N/m. Right: start
  pose, worst direction, 1.25\,mm steps.}
  \label{tab:controller-choices}
  \begin{tabular}{lr@{\qquad}lrr}
    \toprule
    Target                      & travelled & Damping              & following & overshoot \\
    \midrule
    measured position + action  & 1--4\,cm  & Franka, $2\sqrt{K}$   & 1.5\,mm   & 7.1\,mm \\
    previous target + action    & 10\,cm    & apparent mass         & 1.0\,mm   & 2.0\,mm \\
    \bottomrule
  \end{tabular}
\end{table}

\begin{table}[h]
  \centering
  \caption{Largest step: 9 start points over the workspace (tool tip at $x$ 0.35--0.60\,m, $y \pm 0.20$\,m, $z$
  0.20--0.45\,m, and the start pose), 6\,cm moves along $\pm x, \pm y, \pm z$, worst case. Two points close to
  the robot and low ($x$ 0.35\,m, $z$ 0.20\,m) are within 0.2\,rad of the elbow's joint limit and are listed
  separately.}
  \label{tab:controller-steps}
  \begin{tabular}{rrrrrr}
    \toprule
    Step   & Speed      & following & overshoot & tilt        & following near the joint limit \\
    \midrule
    2.0\,mm  & 6.3\,cm/s & 2.25\,mm & 3.9\,mm & 0.96$^\circ$ & 3.5\,mm \\
    1.75\,mm & 5.5\,cm/s & 1.94\,mm & 3.4\,mm & 0.82$^\circ$ & 3.1\,mm \\
    1.5\,mm  & 4.7\,cm/s & 1.63\,mm & 2.9\,mm & 0.68$^\circ$ & 2.7\,mm \\
    \bottomrule
  \end{tabular}

  \bigskip

  \caption{Rotation step: 20$^\circ$ turns about the base axes at the same 9 points, worst case. The tool tip should
  stay in place while the fork turns.}
  \label{tab:controller-rotation}
  \begin{tabular}{rrrrr}
    \toprule
    Step           & Speed               & angle following & angle overshoot & tool-tip drift \\
    \midrule
    0.25$^\circ$ & 7.8$^\circ$/s & 0.7$^\circ$ & 1.0$^\circ$ & 1.7\,mm \\
    0.30$^\circ$ & 9.4$^\circ$/s & 1.1$^\circ$ & 1.1$^\circ$ & 2.1\,mm \\
    1$^\circ$    & 31$^\circ$/s  & 2.9$^\circ$ & 1.5$^\circ$ & 5.7\,mm \\
    \bottomrule
  \end{tabular}
\end{table}

\paragraph{Findings.}
\begin{itemize}
  \item Adding each action to the previous target is essential; with Franka's damping the tool overshoots by up
        to 7\,mm in the directions where the arm feels heavy; damping by the apparent mass removes most of it.
  \item The simulated FR3 initially had no motor inertia. Without it the arm oscillated and pressed the stem with
        5\,N instead of 0.85\,N; at the wrist the motor inertia is about 100 times the link inertia.
  \item Turning the fork moves the tool tip: the law treats position and orientation as independent springs, but
        the arm's dynamics couple them. With the criterion (tip drift $\le 2$\,mm, angles $\le 2^\circ$) the first
        design limited rotations to 7.8$^\circ$/s. The acceleration limit (Section~\ref{sec:al}) now allows
        45$^\circ$/s within the criterion. Decoupling rotation and translation by inertia feedforward
        (operational-space control~\cite{khatib1987}) is option B3 of Section~\ref{sec:al}: planned only after a
        policy trains successfully and the robot's model has been checked (identification), because it makes the
        motion depend on the model.
  \item Close to a joint limit the arm lags more; there the real robot would stop. The task will keep the tool away
        from these poses (workspace bound or penalty).
  \item When pushing the stem, contact force (0.8 to 1.4\,N) and curvature (2.0 to 3.2\,1/m, limit 5) depend on
        the push, hardly on the controller: the stem is far softer than the controller's spring.
\end{itemize}

\paragraph{Limits.} Simulation only; robot parameters from Franka's description files, not identified on our
robot; no joint friction; control period 2\,ms instead of 1\,ms.

\paragraph{Questions.}
\begin{enumerate}
  \item Is this controller acceptable next to a plant ($K_p = 1000$\,N/m, $K_o = 100$\,N\,m/rad, at most 10\,cm/s and
        45$^\circ$/s, contact force bounded to 4\,N)?
  \item Lab setup: is the Franka Control Interface (torque control) enabled on our FR3, with which system and
        libfranka versions; is there a real-time PC for it; does the lab already use a control stack
        (franka\_ros2, frankapy, own C++) that this controller should be added to?
\end{enumerate}

{\raggedright  % long code names in the references: no stretched lines
\begin{thebibliography}{99}
\bibitem{tang2023industreal} B.~Tang, M.~A.~Lin, I.~Akinola, A.~Handa, G.~S.~Sukhatme, F.~Ramos, D.~Fox,
  Y.~Narang. IndustReal: transferring contact-rich assembly tasks from simulation to reality. \emph{RSS}, 2023.
\bibitem{frankatwin} FrankaTwin: aligned simulation and task-space impedance control for Franka robots in Isaac
  Lab. \texttt{github.com/tsrobcvai/frankatwin}.
\bibitem{kaspar2020sim2real} M.~Kaspar, J.~D.~Muñoz Osorio, J.~Bock. Sim2Real transfer for reinforcement learning
  without dynamics randomization. \emph{IEEE/RSJ IROS}, 2020.
\bibitem{peng2018sim} X.~B.~Peng, M.~Andrychowicz, W.~Zaremba, P.~Abbeel. Sim-to-real transfer of robotic control
  with dynamics randomization. \emph{IEEE ICRA}, 2018.
\bibitem{aljalbout2024action} E.~Aljalbout, F.~Frank, M.~Karl, P.~van der Smagt. On the role of the action space
  in robot manipulation learning and sim-to-real transfer. \emph{IEEE Robotics and Automation Letters}, 2024.
\bibitem{aljalbout2026realitygap} E.~Aljalbout, J.~Xing, A.~Romero, I.~Akinola, C.~R.~Garrett, E.~Heiden,
  A.~Gupta, T.~Hermans, Y.~Narang, D.~Fox, D.~Scaramuzza, F.~Ramos. The reality gap in robotics: challenges,
  solutions, and best practices. \emph{Annual Review of Control, Robotics, and Autonomous Systems}, 9, 2026.
\bibitem{albu2003cartesian} A.~Albu-Schäffer, C.~Ott, U.~Frese, G.~Hirzinger. Cartesian impedance control of
  redundant robots: recent results with the DLR-light-weight-arms. \emph{IEEE ICRA}, 2003.
\bibitem{khatib1987} O.~Khatib. A unified approach for motion and force control of robot manipulators: the
  operational space formulation. \emph{IEEE Journal on Robotics and Automation}, 3(1), 1987.
\bibitem{libfranka} Franka Robotics. libfranka example \texttt{cartesian\_impedance\_control.cpp}.
  \texttt{github.com/frankaemika/libfranka}.
\bibitem{frankaros} Franka Robotics. franka\_ros \texttt{cartesian\_impedance\_example\_controller}.
  \texttt{github.com/frankaemika/franka\_ros}.
\end{thebibliography}
}
% ======================================== END CONTENT ========================================
); uses only packages loaded there
% (amsmath, booktabs). Numbers: scripts/sweep_impedance.py, scripts/sweep_action_poses.py,
% scripts/check_push_env.py and docs/notes/2026-10-06.md to 2026-10-08.md in the project repo. Current values and
% their derivation: action_limits_problem.tex (Section sec:al).

% ======================================= BEGIN CONTENT =======================================
\section{Open question: robot controller and action space}
\label{sec:action-controller}

\paragraph{The problem.} The policy does not move the robot's joints directly. At 31\,Hz it outputs a small motion
command for the fork (``move a little this way, turn a little''), and a low-level controller turns these commands
into motor torques at 1\,kHz. The policy learns how the fork reacts to its commands, so two things must hold for it
to work on the real robot: (1) the controller in simulation must be the one that runs on the real FR3, with the same
values, otherwise the fork reacts differently after the transfer; (2) the fork should follow the commanded steps with
a small lag. This is not strictly required, since the policy could learn a lag, but with a small lag the motion is
set by the controller, which is identical in simulation and on the robot, rather than by the arm's dynamics, which
are not; action spaces that keep the tracking error small transfer better~\cite{aljalbout2024action}. The price is
agility (Section~\ref{sec:al}). This section proposes the controller, the form of the commands, and how large a step
may be.

\paragraph{Summary.} The policy commands small steps of the tool's position and orientation. They are executed by
Franka's Cartesian impedance controller, with one change in the damping, and the \emph{same} controller with the
\emph{same} values runs in simulation and on the real FR3 (Section~\ref{sec:controller-what}). Successful sim-to-real
work on Franka robots does the same~\cite{tang2023industreal,frankatwin}. In simulation the tool follows the
commands within 2\,mm (1.68\,mm measured) over the workspace at up to 10\,cm/s and 45$^\circ$/s, with the speed-up
limited to 0.08\,m/s$^2$ and 20$^\circ$/s$^2$ (Section~\ref{sec:al}; the first design, Section~\ref{sec:controller-results},
was limited to 5.5\,cm/s and 7.8$^\circ$/s). We ask the supervisors to
confirm the controller and to tell us about the lab's FR3 setup (end of this section).

\subsection{What we use}
\label{sec:controller-what}

\paragraph{Action.} Every policy step (31.25\,Hz) the policy outputs a displacement of the tool tip (at most
3.2\,mm, i.e.\ 10\,cm/s) and a rotation of the fork (rotation vector, at most 1.44$^\circ$, i.e.\ 45$^\circ$/s), both in
the robot's base frame. The step may change by at most 0.08\,mm (0.02$^\circ$) from one policy step to the next, which
limits the acceleration to 0.08\,m/s$^2$ (20$^\circ$/s$^2$) and with it the lag (Section~\ref{sec:al}). The steps are
applied to the \emph{previous target} pose, not to the measured tool pose (IndustReal's policy-level action
integrator~\cite{tang2023industreal}), so commanded steps are not lost while the tool moves freely. Within the policy
step the target moves and turns at constant rate, and its velocities $\dot x_d$, $\omega_d$ are passed to the
controller. At every control step the target is kept within 4\,mm and 3$^\circ$ of the measured tool pose, which bounds
the spring force to 4\,N and the moment to 5.2\,N\,m when the fork is blocked; while held, the target moves with the
tool and its velocity is the tool's. The fork starts horizontal; its orientation is then free.

\paragraph{Control law.} Every control step the joint torques are
\begin{equation}
  \tau = J^\top \begin{bmatrix} K_p\,(x_d - x) - D_p\,(\dot x - \dot x_d) \\
                                 -K_o\,e_o - D_o\,(\omega - \omega_d) \end{bmatrix}
       + \left(I - J^\top J^{\top+}\right)\left(k_\text{ns}(q_\text{ns} - q) - 2\sqrt{k_\text{ns}}\,\dot q\right)
       + \tau_\text{Coriolis},
  \label{eq:impedance}
\end{equation}
a spring and damper that pull the tool tip $x$ to its target $x_d$ and its orientation to the target orientation
($e_o$: orientation error), plus a weak spring that keeps the elbow of the redundant arm near its rest posture. This is the law of
Franka's example controllers~\cite{libfranka,frankaros}, with one change: Franka sets the damping to
$D = 2\sqrt{K}$, which is critically damped only for a tool of 1\,kg. The arm's inertia felt at the tool,
$\Lambda = (J M^{-1} J^\top)^{-1}$, is several kilograms and depends on pose and direction. We therefore use
$D = \Lambda^{1/2} K^{1/2} + K^{1/2} \Lambda^{1/2}$~\cite{albu2003cartesian}, i.e.\ $2\sqrt{K}\,\Lambda^{1/2}$,
critically damped everywhere. Values: $K_p = 1000$\,N/m, $K_o = 100$\,N\,m/rad (stiffer than Franka's examples,
which use up to 400\,N/m and 30\,N\,m/rad, to keep the lag small; Section~\ref{sec:al}), $k_\text{ns} = 0.5$\,N\,m/rad,
torque change at most 1\,N\,m per ms (Franka's limit).

\paragraph{Simulation and real robot.}
\begin{center}
\begin{tabular}{p{0.2\textwidth}p{0.36\textwidth}p{0.36\textwidth}}
  \toprule
  & Simulation (Isaac Lab, PhysX) & Real FR3 \\
  \midrule
  Law~\eqref{eq:impedance} & Python function, every 2\,ms, joint torques to the simulator
    & the same equations in C++ in a libfranka torque loop at 1\,kHz \\
  $J$, $M$, Coriolis & from the simulator (Coriolis left out; small below 0.1\,m/s)
    & from Franka's robot model in libfranka \\
  Gravity & switched off for the robot & compensated by the robot (fork mass entered as load) \\
  Motor inertia & added to every joint (gear ratio$^2\times$motor inertia, from Franka's description files)
    & physical; whether libfranka's $M$ includes it is to be checked \\
  Target clamp, feedforward & every 2\,ms, function \texttt{substep\_command} & the same function in C++ in the 1\,kHz loop \\
  Speed cap, acceleration limit & policy side, every 32\,ms & the same Python code, every 32\,ms \\
  Policy & same process, every 32\,ms & Python process, sends the step to the C++ loop every 32\,ms \\
  \bottomrule
\end{tabular}
\end{center}
To make sure both sides compute the same torques, the C++ version will be checked against the Python version on
recorded robot states. Remaining differences (friction, exact inertia, delays) are handled by measuring the real
robot's response to the same commanded paths and by randomizing gains and timing during
training~\cite{peng2018sim,frankatwin}.

\subsection{Why}

The literature that transferred contact-rich tasks to a real Franka runs the same task-space impedance
controller with the same gains in simulation and on the robot~\cite{tang2023industreal,frankatwin,kaspar2020sim2real},
and identifies or randomizes what still differs~\cite{peng2018sim}. The choice of action space measurably affects
transfer~\cite{aljalbout2024action}; for an overview of the reality gap see~\cite{aljalbout2026realitygap}.
The placeholder we started from (inverse kinematics and stiff joint position control with Isaac Lab's default gains)
lagged 0.2\,s behind its targets and has no counterpart on the real robot.

\subsection{Results in simulation}
\label{sec:controller-results}

\emph{These results led to the first design (speed cap only, $K_p = 400$\,N/m, $K_o = 30$\,N\,m/rad, at most
5.5\,cm/s). They are kept as history; the current values and measurements are in Section~\ref{sec:al}.}

Criterion for ``the step is followed'': while moving, the tool trails the commanded path by at most 2\,mm, and it
overshoots the end point by at most 5\,mm (fork tilt at most 1$^\circ$); while turning, the angle lags and overshoots
by at most 2$^\circ$ and the tool tip moves by at most 2\,mm. Every step is eventually executed: after a
hold the tool is within 0.5\,mm of the commanded end point in all tests.

\begin{table}[h]
  \centering
  \caption{Choices tested (fork in free space). Left: 10\,cm commanded, 5\,cm/s, $K_p = 150$\,N/m. Right: start
  pose, worst direction, 1.25\,mm steps.}
  \label{tab:controller-choices}
  \begin{tabular}{lr@{\qquad}lrr}
    \toprule
    Target                      & travelled & Damping              & following & overshoot \\
    \midrule
    measured position + action  & 1--4\,cm  & Franka, $2\sqrt{K}$   & 1.5\,mm   & 7.1\,mm \\
    previous target + action    & 10\,cm    & apparent mass         & 1.0\,mm   & 2.0\,mm \\
    \bottomrule
  \end{tabular}
\end{table}

\begin{table}[h]
  \centering
  \caption{Largest step: 9 start points over the workspace (tool tip at $x$ 0.35--0.60\,m, $y \pm 0.20$\,m, $z$
  0.20--0.45\,m, and the start pose), 6\,cm moves along $\pm x, \pm y, \pm z$, worst case. Two points close to
  the robot and low ($x$ 0.35\,m, $z$ 0.20\,m) are within 0.2\,rad of the elbow's joint limit and are listed
  separately.}
  \label{tab:controller-steps}
  \begin{tabular}{rrrrrr}
    \toprule
    Step   & Speed      & following & overshoot & tilt        & following near the joint limit \\
    \midrule
    2.0\,mm  & 6.3\,cm/s & 2.25\,mm & 3.9\,mm & 0.96$^\circ$ & 3.5\,mm \\
    1.75\,mm & 5.5\,cm/s & 1.94\,mm & 3.4\,mm & 0.82$^\circ$ & 3.1\,mm \\
    1.5\,mm  & 4.7\,cm/s & 1.63\,mm & 2.9\,mm & 0.68$^\circ$ & 2.7\,mm \\
    \bottomrule
  \end{tabular}

  \bigskip

  \caption{Rotation step: 20$^\circ$ turns about the base axes at the same 9 points, worst case. The tool tip should
  stay in place while the fork turns.}
  \label{tab:controller-rotation}
  \begin{tabular}{rrrrr}
    \toprule
    Step           & Speed               & angle following & angle overshoot & tool-tip drift \\
    \midrule
    0.25$^\circ$ & 7.8$^\circ$/s & 0.7$^\circ$ & 1.0$^\circ$ & 1.7\,mm \\
    0.30$^\circ$ & 9.4$^\circ$/s & 1.1$^\circ$ & 1.1$^\circ$ & 2.1\,mm \\
    1$^\circ$    & 31$^\circ$/s  & 2.9$^\circ$ & 1.5$^\circ$ & 5.7\,mm \\
    \bottomrule
  \end{tabular}
\end{table}

\paragraph{Findings.}
\begin{itemize}
  \item Adding each action to the previous target is essential; with Franka's damping the tool overshoots by up
        to 7\,mm in the directions where the arm feels heavy; damping by the apparent mass removes most of it.
  \item The simulated FR3 initially had no motor inertia. Without it the arm oscillated and pressed the stem with
        5\,N instead of 0.85\,N; at the wrist the motor inertia is about 100 times the link inertia.
  \item Turning the fork moves the tool tip: the law treats position and orientation as independent springs, but
        the arm's dynamics couple them. With the criterion (tip drift $\le 2$\,mm, angles $\le 2^\circ$) the first
        design limited rotations to 7.8$^\circ$/s. The acceleration limit (Section~\ref{sec:al}) now allows
        45$^\circ$/s within the criterion. Decoupling rotation and translation by inertia feedforward
        (operational-space control~\cite{khatib1987}) is option B3 of Section~\ref{sec:al}: planned only after a
        policy trains successfully and the robot's model has been checked (identification), because it makes the
        motion depend on the model.
  \item Close to a joint limit the arm lags more; there the real robot would stop. The task will keep the tool away
        from these poses (workspace bound or penalty).
  \item When pushing the stem, contact force (0.8 to 1.4\,N) and curvature (2.0 to 3.2\,1/m, limit 5) depend on
        the push, hardly on the controller: the stem is far softer than the controller's spring.
\end{itemize}

\paragraph{Limits.} Simulation only; robot parameters from Franka's description files, not identified on our
robot; no joint friction; control period 2\,ms instead of 1\,ms.

\paragraph{Questions.}
\begin{enumerate}
  \item Is this controller acceptable next to a plant ($K_p = 1000$\,N/m, $K_o = 100$\,N\,m/rad, at most 10\,cm/s and
        45$^\circ$/s, contact force bounded to 4\,N)?
  \item Lab setup: is the Franka Control Interface (torque control) enabled on our FR3, with which system and
        libfranka versions; is there a real-time PC for it; does the lab already use a control stack
        (franka\_ros2, frankapy, own C++) that this controller should be added to?
\end{enumerate}

{\raggedright  % long code names in the references: no stretched lines
\begin{thebibliography}{99}
\bibitem{tang2023industreal} B.~Tang, M.~A.~Lin, I.~Akinola, A.~Handa, G.~S.~Sukhatme, F.~Ramos, D.~Fox,
  Y.~Narang. IndustReal: transferring contact-rich assembly tasks from simulation to reality. \emph{RSS}, 2023.
\bibitem{frankatwin} FrankaTwin: aligned simulation and task-space impedance control for Franka robots in Isaac
  Lab. \texttt{github.com/tsrobcvai/frankatwin}.
\bibitem{kaspar2020sim2real} M.~Kaspar, J.~D.~Muñoz Osorio, J.~Bock. Sim2Real transfer for reinforcement learning
  without dynamics randomization. \emph{IEEE/RSJ IROS}, 2020.
\bibitem{peng2018sim} X.~B.~Peng, M.~Andrychowicz, W.~Zaremba, P.~Abbeel. Sim-to-real transfer of robotic control
  with dynamics randomization. \emph{IEEE ICRA}, 2018.
\bibitem{aljalbout2024action} E.~Aljalbout, F.~Frank, M.~Karl, P.~van der Smagt. On the role of the action space
  in robot manipulation learning and sim-to-real transfer. \emph{IEEE Robotics and Automation Letters}, 2024.
\bibitem{aljalbout2026realitygap} E.~Aljalbout, J.~Xing, A.~Romero, I.~Akinola, C.~R.~Garrett, E.~Heiden,
  A.~Gupta, T.~Hermans, Y.~Narang, D.~Fox, D.~Scaramuzza, F.~Ramos. The reality gap in robotics: challenges,
  solutions, and best practices. \emph{Annual Review of Control, Robotics, and Autonomous Systems}, 9, 2026.
\bibitem{albu2003cartesian} A.~Albu-Schäffer, C.~Ott, U.~Frese, G.~Hirzinger. Cartesian impedance control of
  redundant robots: recent results with the DLR-light-weight-arms. \emph{IEEE ICRA}, 2003.
\bibitem{khatib1987} O.~Khatib. A unified approach for motion and force control of robot manipulators: the
  operational space formulation. \emph{IEEE Journal on Robotics and Automation}, 3(1), 1987.
\bibitem{libfranka} Franka Robotics. libfranka example \texttt{cartesian\_impedance\_control.cpp}.
  \texttt{github.com/frankaemika/libfranka}.
\bibitem{frankaros} Franka Robotics. franka\_ros \texttt{cartesian\_impedance\_example\_controller}.
  \texttt{github.com/frankaemika/franka\_ros}.
\end{thebibliography}
}
% ======================================== END CONTENT ========================================
); uses only packages loaded there
% (amsmath, booktabs). Numbers: scripts/sweep_impedance.py, scripts/sweep_action_poses.py,
% scripts/check_push_env.py and docs/notes/2026-10-06.md to 2026-10-08.md in the project repo. Current values and
% their derivation: action_limits_problem.tex (Section sec:al).

% ======================================= BEGIN CONTENT =======================================
\section{Open question: robot controller and action space}
\label{sec:action-controller}

\paragraph{The problem.} The policy does not move the robot's joints directly. At 31\,Hz it outputs a small motion
command for the fork (``move a little this way, turn a little''), and a low-level controller turns these commands
into motor torques at 1\,kHz. The policy learns how the fork reacts to its commands, so two things must hold for it
to work on the real robot: (1) the controller in simulation must be the one that runs on the real FR3, with the same
values, otherwise the fork reacts differently after the transfer; (2) the fork should follow the commanded steps with
a small lag. This is not strictly required, since the policy could learn a lag, but with a small lag the motion is
set by the controller, which is identical in simulation and on the robot, rather than by the arm's dynamics, which
are not; action spaces that keep the tracking error small transfer better~\cite{aljalbout2024action}. The price is
agility (Section~\ref{sec:al}). This section proposes the controller, the form of the commands, and how large a step
may be.

\paragraph{Summary.} The policy commands small steps of the tool's position and orientation. They are executed by
Franka's Cartesian impedance controller, with one change in the damping, and the \emph{same} controller with the
\emph{same} values runs in simulation and on the real FR3 (Section~\ref{sec:controller-what}). Successful sim-to-real
work on Franka robots does the same~\cite{tang2023industreal,frankatwin}. In simulation the tool follows the
commands within 2\,mm (1.68\,mm measured) over the workspace at up to 10\,cm/s and 45$^\circ$/s, with the speed-up
limited to 0.08\,m/s$^2$ and 20$^\circ$/s$^2$ (Section~\ref{sec:al}; the first design, Section~\ref{sec:controller-results},
was limited to 5.5\,cm/s and 7.8$^\circ$/s). We ask the supervisors to
confirm the controller and to tell us about the lab's FR3 setup (end of this section).

\subsection{What we use}
\label{sec:controller-what}

\paragraph{Action.} Every policy step (31.25\,Hz) the policy outputs a displacement of the tool tip (at most
3.2\,mm, i.e.\ 10\,cm/s) and a rotation of the fork (rotation vector, at most 1.44$^\circ$, i.e.\ 45$^\circ$/s), both in
the robot's base frame. The step may change by at most 0.08\,mm (0.02$^\circ$) from one policy step to the next, which
limits the acceleration to 0.08\,m/s$^2$ (20$^\circ$/s$^2$) and with it the lag (Section~\ref{sec:al}). The steps are
applied to the \emph{previous target} pose, not to the measured tool pose (IndustReal's policy-level action
integrator~\cite{tang2023industreal}), so commanded steps are not lost while the tool moves freely. Within the policy
step the target moves and turns at constant rate, and its velocities $\dot x_d$, $\omega_d$ are passed to the
controller. At every control step the target is kept within 4\,mm and 3$^\circ$ of the measured tool pose, which bounds
the spring force to 4\,N and the moment to 5.2\,N\,m when the fork is blocked; while held, the target moves with the
tool and its velocity is the tool's. The fork starts horizontal; its orientation is then free.

\paragraph{Control law.} Every control step the joint torques are
\begin{equation}
  \tau = J^\top \begin{bmatrix} K_p\,(x_d - x) - D_p\,(\dot x - \dot x_d) \\
                                 -K_o\,e_o - D_o\,(\omega - \omega_d) \end{bmatrix}
       + \left(I - J^\top J^{\top+}\right)\left(k_\text{ns}(q_\text{ns} - q) - 2\sqrt{k_\text{ns}}\,\dot q\right)
       + \tau_\text{Coriolis},
  \label{eq:impedance}
\end{equation}
a spring and damper that pull the tool tip $x$ to its target $x_d$ and its orientation to the target orientation
($e_o$: orientation error), plus a weak spring that keeps the elbow of the redundant arm near its rest posture. This is the law of
Franka's example controllers~\cite{libfranka,frankaros}, with one change: Franka sets the damping to
$D = 2\sqrt{K}$, which is critically damped only for a tool of 1\,kg. The arm's inertia felt at the tool,
$\Lambda = (J M^{-1} J^\top)^{-1}$, is several kilograms and depends on pose and direction. We therefore use
$D = \Lambda^{1/2} K^{1/2} + K^{1/2} \Lambda^{1/2}$~\cite{albu2003cartesian}, i.e.\ $2\sqrt{K}\,\Lambda^{1/2}$,
critically damped everywhere. Values: $K_p = 1000$\,N/m, $K_o = 100$\,N\,m/rad (stiffer than Franka's examples,
which use up to 400\,N/m and 30\,N\,m/rad, to keep the lag small; Section~\ref{sec:al}), $k_\text{ns} = 0.5$\,N\,m/rad,
torque change at most 1\,N\,m per ms (Franka's limit).

\paragraph{Simulation and real robot.}
\begin{center}
\begin{tabular}{p{0.2\textwidth}p{0.36\textwidth}p{0.36\textwidth}}
  \toprule
  & Simulation (Isaac Lab, PhysX) & Real FR3 \\
  \midrule
  Law~\eqref{eq:impedance} & Python function, every 2\,ms, joint torques to the simulator
    & the same equations in C++ in a libfranka torque loop at 1\,kHz \\
  $J$, $M$, Coriolis & from the simulator (Coriolis left out; small below 0.1\,m/s)
    & from Franka's robot model in libfranka \\
  Gravity & switched off for the robot & compensated by the robot (fork mass entered as load) \\
  Motor inertia & added to every joint (gear ratio$^2\times$motor inertia, from Franka's description files)
    & physical; whether libfranka's $M$ includes it is to be checked \\
  Target clamp, feedforward & every 2\,ms, function \texttt{substep\_command} & the same function in C++ in the 1\,kHz loop \\
  Speed cap, acceleration limit & policy side, every 32\,ms & the same Python code, every 32\,ms \\
  Policy & same process, every 32\,ms & Python process, sends the step to the C++ loop every 32\,ms \\
  \bottomrule
\end{tabular}
\end{center}
To make sure both sides compute the same torques, the C++ version will be checked against the Python version on
recorded robot states. Remaining differences (friction, exact inertia, delays) are handled by measuring the real
robot's response to the same commanded paths and by randomizing gains and timing during
training~\cite{peng2018sim,frankatwin}.

\subsection{Why}

The literature that transferred contact-rich tasks to a real Franka runs the same task-space impedance
controller with the same gains in simulation and on the robot~\cite{tang2023industreal,frankatwin,kaspar2020sim2real},
and identifies or randomizes what still differs~\cite{peng2018sim}. The choice of action space measurably affects
transfer~\cite{aljalbout2024action}; for an overview of the reality gap see~\cite{aljalbout2026realitygap}.
The placeholder we started from (inverse kinematics and stiff joint position control with Isaac Lab's default gains)
lagged 0.2\,s behind its targets and has no counterpart on the real robot.

\subsection{Results in simulation}
\label{sec:controller-results}

\emph{These results led to the first design (speed cap only, $K_p = 400$\,N/m, $K_o = 30$\,N\,m/rad, at most
5.5\,cm/s). They are kept as history; the current values and measurements are in Section~\ref{sec:al}.}

Criterion for ``the step is followed'': while moving, the tool trails the commanded path by at most 2\,mm, and it
overshoots the end point by at most 5\,mm (fork tilt at most 1$^\circ$); while turning, the angle lags and overshoots
by at most 2$^\circ$ and the tool tip moves by at most 2\,mm. Every step is eventually executed: after a
hold the tool is within 0.5\,mm of the commanded end point in all tests.

\begin{table}[h]
  \centering
  \caption{Choices tested (fork in free space). Left: 10\,cm commanded, 5\,cm/s, $K_p = 150$\,N/m. Right: start
  pose, worst direction, 1.25\,mm steps.}
  \label{tab:controller-choices}
  \begin{tabular}{lr@{\qquad}lrr}
    \toprule
    Target                      & travelled & Damping              & following & overshoot \\
    \midrule
    measured position + action  & 1--4\,cm  & Franka, $2\sqrt{K}$   & 1.5\,mm   & 7.1\,mm \\
    previous target + action    & 10\,cm    & apparent mass         & 1.0\,mm   & 2.0\,mm \\
    \bottomrule
  \end{tabular}
\end{table}

\begin{table}[h]
  \centering
  \caption{Largest step: 9 start points over the workspace (tool tip at $x$ 0.35--0.60\,m, $y \pm 0.20$\,m, $z$
  0.20--0.45\,m, and the start pose), 6\,cm moves along $\pm x, \pm y, \pm z$, worst case. Two points close to
  the robot and low ($x$ 0.35\,m, $z$ 0.20\,m) are within 0.2\,rad of the elbow's joint limit and are listed
  separately.}
  \label{tab:controller-steps}
  \begin{tabular}{rrrrrr}
    \toprule
    Step   & Speed      & following & overshoot & tilt        & following near the joint limit \\
    \midrule
    2.0\,mm  & 6.3\,cm/s & 2.25\,mm & 3.9\,mm & 0.96$^\circ$ & 3.5\,mm \\
    1.75\,mm & 5.5\,cm/s & 1.94\,mm & 3.4\,mm & 0.82$^\circ$ & 3.1\,mm \\
    1.5\,mm  & 4.7\,cm/s & 1.63\,mm & 2.9\,mm & 0.68$^\circ$ & 2.7\,mm \\
    \bottomrule
  \end{tabular}

  \bigskip

  \caption{Rotation step: 20$^\circ$ turns about the base axes at the same 9 points, worst case. The tool tip should
  stay in place while the fork turns.}
  \label{tab:controller-rotation}
  \begin{tabular}{rrrrr}
    \toprule
    Step           & Speed               & angle following & angle overshoot & tool-tip drift \\
    \midrule
    0.25$^\circ$ & 7.8$^\circ$/s & 0.7$^\circ$ & 1.0$^\circ$ & 1.7\,mm \\
    0.30$^\circ$ & 9.4$^\circ$/s & 1.1$^\circ$ & 1.1$^\circ$ & 2.1\,mm \\
    1$^\circ$    & 31$^\circ$/s  & 2.9$^\circ$ & 1.5$^\circ$ & 5.7\,mm \\
    \bottomrule
  \end{tabular}
\end{table}

\paragraph{Findings.}
\begin{itemize}
  \item Adding each action to the previous target is essential; with Franka's damping the tool overshoots by up
        to 7\,mm in the directions where the arm feels heavy; damping by the apparent mass removes most of it.
  \item The simulated FR3 initially had no motor inertia. Without it the arm oscillated and pressed the stem with
        5\,N instead of 0.85\,N; at the wrist the motor inertia is about 100 times the link inertia.
  \item Turning the fork moves the tool tip: the law treats position and orientation as independent springs, but
        the arm's dynamics couple them. With the criterion (tip drift $\le 2$\,mm, angles $\le 2^\circ$) the first
        design limited rotations to 7.8$^\circ$/s. The acceleration limit (Section~\ref{sec:al}) now allows
        45$^\circ$/s within the criterion. Decoupling rotation and translation by inertia feedforward
        (operational-space control~\cite{khatib1987}) is option B3 of Section~\ref{sec:al}: planned only after a
        policy trains successfully and the robot's model has been checked (identification), because it makes the
        motion depend on the model.
  \item Close to a joint limit the arm lags more; there the real robot would stop. The task will keep the tool away
        from these poses (workspace bound or penalty).
  \item When pushing the stem, contact force (0.8 to 1.4\,N) and curvature (2.0 to 3.2\,1/m, limit 5) depend on
        the push, hardly on the controller: the stem is far softer than the controller's spring.
\end{itemize}

\paragraph{Limits.} Simulation only; robot parameters from Franka's description files, not identified on our
robot; no joint friction; control period 2\,ms instead of 1\,ms.

\paragraph{Questions.}
\begin{enumerate}
  \item Is this controller acceptable next to a plant ($K_p = 1000$\,N/m, $K_o = 100$\,N\,m/rad, at most 10\,cm/s and
        45$^\circ$/s, contact force bounded to 4\,N)?
  \item Lab setup: is the Franka Control Interface (torque control) enabled on our FR3, with which system and
        libfranka versions; is there a real-time PC for it; does the lab already use a control stack
        (franka\_ros2, frankapy, own C++) that this controller should be added to?
\end{enumerate}

{\raggedright  % long code names in the references: no stretched lines
\begin{thebibliography}{99}
\bibitem{tang2023industreal} B.~Tang, M.~A.~Lin, I.~Akinola, A.~Handa, G.~S.~Sukhatme, F.~Ramos, D.~Fox,
  Y.~Narang. IndustReal: transferring contact-rich assembly tasks from simulation to reality. \emph{RSS}, 2023.
\bibitem{frankatwin} FrankaTwin: aligned simulation and task-space impedance control for Franka robots in Isaac
  Lab. \texttt{github.com/tsrobcvai/frankatwin}.
\bibitem{kaspar2020sim2real} M.~Kaspar, J.~D.~Muñoz Osorio, J.~Bock. Sim2Real transfer for reinforcement learning
  without dynamics randomization. \emph{IEEE/RSJ IROS}, 2020.
\bibitem{peng2018sim} X.~B.~Peng, M.~Andrychowicz, W.~Zaremba, P.~Abbeel. Sim-to-real transfer of robotic control
  with dynamics randomization. \emph{IEEE ICRA}, 2018.
\bibitem{aljalbout2024action} E.~Aljalbout, F.~Frank, M.~Karl, P.~van der Smagt. On the role of the action space
  in robot manipulation learning and sim-to-real transfer. \emph{IEEE Robotics and Automation Letters}, 2024.
\bibitem{aljalbout2026realitygap} E.~Aljalbout, J.~Xing, A.~Romero, I.~Akinola, C.~R.~Garrett, E.~Heiden,
  A.~Gupta, T.~Hermans, Y.~Narang, D.~Fox, D.~Scaramuzza, F.~Ramos. The reality gap in robotics: challenges,
  solutions, and best practices. \emph{Annual Review of Control, Robotics, and Autonomous Systems}, 9, 2026.
\bibitem{albu2003cartesian} A.~Albu-Schäffer, C.~Ott, U.~Frese, G.~Hirzinger. Cartesian impedance control of
  redundant robots: recent results with the DLR-light-weight-arms. \emph{IEEE ICRA}, 2003.
\bibitem{khatib1987} O.~Khatib. A unified approach for motion and force control of robot manipulators: the
  operational space formulation. \emph{IEEE Journal on Robotics and Automation}, 3(1), 1987.
\bibitem{libfranka} Franka Robotics. libfranka example \texttt{cartesian\_impedance\_control.cpp}.
  \texttt{github.com/frankaemika/libfranka}.
\bibitem{frankaros} Franka Robotics. franka\_ros \texttt{cartesian\_impedance\_example\_controller}.
  \texttt{github.com/frankaemika/franka\_ros}.
\end{thebibliography}
}
% ======================================== END CONTENT ========================================
); uses only packages loaded there
% (amsmath, booktabs). Numbers: scripts/sweep_impedance.py, scripts/sweep_action_poses.py,
% scripts/check_push_env.py and docs/notes/2026-10-06.md to 2026-10-08.md in the project repo. Current values and
% their derivation: action_limits_problem.tex (Section sec:al).

% ======================================= BEGIN CONTENT =======================================
\section{Open question: robot controller and action space}
\label{sec:action-controller}

\paragraph{The problem.} The policy does not move the robot's joints directly. At 31\,Hz it outputs a small motion
command for the fork (``move a little this way, turn a little''), and a low-level controller turns these commands
into motor torques at 1\,kHz. The policy learns how the fork reacts to its commands, so two things must hold for it
to work on the real robot: (1) the controller in simulation must be the one that runs on the real FR3, with the same
values, otherwise the fork reacts differently after the transfer; (2) the fork should follow the commanded steps with
a small lag. This is not strictly required, since the policy could learn a lag, but with a small lag the motion is
set by the controller, which is identical in simulation and on the robot, rather than by the arm's dynamics, which
are not; action spaces that keep the tracking error small transfer better~\cite{aljalbout2024action}. The price is
agility (Section~\ref{sec:al}). This section proposes the controller, the form of the commands, and how large a step
may be.

\paragraph{Summary.} The policy commands small steps of the tool's position and orientation. They are executed by
Franka's Cartesian impedance controller, with one change in the damping, and the \emph{same} controller with the
\emph{same} values runs in simulation and on the real FR3 (Section~\ref{sec:controller-what}). Successful sim-to-real
work on Franka robots does the same~\cite{tang2023industreal,frankatwin}. In simulation the tool follows the
commands within 2\,mm (1.68\,mm measured) over the workspace at up to 10\,cm/s and 45$^\circ$/s, with the speed-up
limited to 0.08\,m/s$^2$ and 20$^\circ$/s$^2$ (Section~\ref{sec:al}; the first design, Section~\ref{sec:controller-results},
was limited to 5.5\,cm/s and 7.8$^\circ$/s). We ask the supervisors to
confirm the controller and to tell us about the lab's FR3 setup (end of this section).

\subsection{What we use}
\label{sec:controller-what}

\paragraph{Action.} Every policy step (31.25\,Hz) the policy outputs a displacement of the tool tip (at most
3.2\,mm, i.e.\ 10\,cm/s) and a rotation of the fork (rotation vector, at most 1.44$^\circ$, i.e.\ 45$^\circ$/s), both in
the robot's base frame. The step may change by at most 0.08\,mm (0.02$^\circ$) from one policy step to the next, which
limits the acceleration to 0.08\,m/s$^2$ (20$^\circ$/s$^2$) and with it the lag (Section~\ref{sec:al}). The steps are
applied to the \emph{previous target} pose, not to the measured tool pose (IndustReal's policy-level action
integrator~\cite{tang2023industreal}), so commanded steps are not lost while the tool moves freely. Within the policy
step the target moves and turns at constant rate, and its velocities $\dot x_d$, $\omega_d$ are passed to the
controller. At every control step the target is kept within 4\,mm and 3$^\circ$ of the measured tool pose, which bounds
the spring force to 4\,N and the moment to 5.2\,N\,m when the fork is blocked; while held, the target moves with the
tool and its velocity is the tool's. The fork starts horizontal; its orientation is then free.

\paragraph{Control law.} Every control step the joint torques are
\begin{equation}
  \tau = J^\top \begin{bmatrix} K_p\,(x_d - x) - D_p\,(\dot x - \dot x_d) \\
                                 -K_o\,e_o - D_o\,(\omega - \omega_d) \end{bmatrix}
       + \left(I - J^\top J^{\top+}\right)\left(k_\text{ns}(q_\text{ns} - q) - 2\sqrt{k_\text{ns}}\,\dot q\right)
       + \tau_\text{Coriolis},
  \label{eq:impedance}
\end{equation}
a spring and damper that pull the tool tip $x$ to its target $x_d$ and its orientation to the target orientation
($e_o$: orientation error), plus a weak spring that keeps the elbow of the redundant arm near its rest posture. This is the law of
Franka's example controllers~\cite{libfranka,frankaros}, with one change: Franka sets the damping to
$D = 2\sqrt{K}$, which is critically damped only for a tool of 1\,kg. The arm's inertia felt at the tool,
$\Lambda = (J M^{-1} J^\top)^{-1}$, is several kilograms and depends on pose and direction. We therefore use
$D = \Lambda^{1/2} K^{1/2} + K^{1/2} \Lambda^{1/2}$~\cite{albu2003cartesian}, i.e.\ $2\sqrt{K}\,\Lambda^{1/2}$,
critically damped everywhere. Values: $K_p = 1000$\,N/m, $K_o = 100$\,N\,m/rad (stiffer than Franka's examples,
which use up to 400\,N/m and 30\,N\,m/rad, to keep the lag small; Section~\ref{sec:al}), $k_\text{ns} = 0.5$\,N\,m/rad,
torque change at most 1\,N\,m per ms (Franka's limit).

\paragraph{Simulation and real robot.}
\begin{center}
\begin{tabular}{p{0.2\textwidth}p{0.36\textwidth}p{0.36\textwidth}}
  \toprule
  & Simulation (Isaac Lab, PhysX) & Real FR3 \\
  \midrule
  Law~\eqref{eq:impedance} & Python function, every 2\,ms, joint torques to the simulator
    & the same equations in C++ in a libfranka torque loop at 1\,kHz \\
  $J$, $M$, Coriolis & from the simulator (Coriolis left out; small below 0.1\,m/s)
    & from Franka's robot model in libfranka \\
  Gravity & switched off for the robot & compensated by the robot (fork mass entered as load) \\
  Motor inertia & added to every joint (gear ratio$^2\times$motor inertia, from Franka's description files)
    & physical; whether libfranka's $M$ includes it is to be checked \\
  Target clamp, feedforward & every 2\,ms, function \texttt{substep\_command} & the same function in C++ in the 1\,kHz loop \\
  Speed cap, acceleration limit & policy side, every 32\,ms & the same Python code, every 32\,ms \\
  Policy & same process, every 32\,ms & Python process, sends the step to the C++ loop every 32\,ms \\
  \bottomrule
\end{tabular}
\end{center}
To make sure both sides compute the same torques, the C++ version will be checked against the Python version on
recorded robot states. Remaining differences (friction, exact inertia, delays) are handled by measuring the real
robot's response to the same commanded paths and by randomizing gains and timing during
training~\cite{peng2018sim,frankatwin}.

\subsection{Why}

The literature that transferred contact-rich tasks to a real Franka runs the same task-space impedance
controller with the same gains in simulation and on the robot~\cite{tang2023industreal,frankatwin,kaspar2020sim2real},
and identifies or randomizes what still differs~\cite{peng2018sim}. The choice of action space measurably affects
transfer~\cite{aljalbout2024action}; for an overview of the reality gap see~\cite{aljalbout2026realitygap}.
The placeholder we started from (inverse kinematics and stiff joint position control with Isaac Lab's default gains)
lagged 0.2\,s behind its targets and has no counterpart on the real robot.

\subsection{Results in simulation}
\label{sec:controller-results}

\emph{These results led to the first design (speed cap only, $K_p = 400$\,N/m, $K_o = 30$\,N\,m/rad, at most
5.5\,cm/s). They are kept as history; the current values and measurements are in Section~\ref{sec:al}.}

Criterion for ``the step is followed'': while moving, the tool trails the commanded path by at most 2\,mm, and it
overshoots the end point by at most 5\,mm (fork tilt at most 1$^\circ$); while turning, the angle lags and overshoots
by at most 2$^\circ$ and the tool tip moves by at most 2\,mm. Every step is eventually executed: after a
hold the tool is within 0.5\,mm of the commanded end point in all tests.

\begin{table}[h]
  \centering
  \caption{Choices tested (fork in free space). Left: 10\,cm commanded, 5\,cm/s, $K_p = 150$\,N/m. Right: start
  pose, worst direction, 1.25\,mm steps.}
  \label{tab:controller-choices}
  \begin{tabular}{lr@{\qquad}lrr}
    \toprule
    Target                      & travelled & Damping              & following & overshoot \\
    \midrule
    measured position + action  & 1--4\,cm  & Franka, $2\sqrt{K}$   & 1.5\,mm   & 7.1\,mm \\
    previous target + action    & 10\,cm    & apparent mass         & 1.0\,mm   & 2.0\,mm \\
    \bottomrule
  \end{tabular}
\end{table}

\begin{table}[h]
  \centering
  \caption{Largest step: 9 start points over the workspace (tool tip at $x$ 0.35--0.60\,m, $y \pm 0.20$\,m, $z$
  0.20--0.45\,m, and the start pose), 6\,cm moves along $\pm x, \pm y, \pm z$, worst case. Two points close to
  the robot and low ($x$ 0.35\,m, $z$ 0.20\,m) are within 0.2\,rad of the elbow's joint limit and are listed
  separately.}
  \label{tab:controller-steps}
  \begin{tabular}{rrrrrr}
    \toprule
    Step   & Speed      & following & overshoot & tilt        & following near the joint limit \\
    \midrule
    2.0\,mm  & 6.3\,cm/s & 2.25\,mm & 3.9\,mm & 0.96$^\circ$ & 3.5\,mm \\
    1.75\,mm & 5.5\,cm/s & 1.94\,mm & 3.4\,mm & 0.82$^\circ$ & 3.1\,mm \\
    1.5\,mm  & 4.7\,cm/s & 1.63\,mm & 2.9\,mm & 0.68$^\circ$ & 2.7\,mm \\
    \bottomrule
  \end{tabular}

  \bigskip

  \caption{Rotation step: 20$^\circ$ turns about the base axes at the same 9 points, worst case. The tool tip should
  stay in place while the fork turns.}
  \label{tab:controller-rotation}
  \begin{tabular}{rrrrr}
    \toprule
    Step           & Speed               & angle following & angle overshoot & tool-tip drift \\
    \midrule
    0.25$^\circ$ & 7.8$^\circ$/s & 0.7$^\circ$ & 1.0$^\circ$ & 1.7\,mm \\
    0.30$^\circ$ & 9.4$^\circ$/s & 1.1$^\circ$ & 1.1$^\circ$ & 2.1\,mm \\
    1$^\circ$    & 31$^\circ$/s  & 2.9$^\circ$ & 1.5$^\circ$ & 5.7\,mm \\
    \bottomrule
  \end{tabular}
\end{table}

\paragraph{Findings.}
\begin{itemize}
  \item Adding each action to the previous target is essential; with Franka's damping the tool overshoots by up
        to 7\,mm in the directions where the arm feels heavy; damping by the apparent mass removes most of it.
  \item The simulated FR3 initially had no motor inertia. Without it the arm oscillated and pressed the stem with
        5\,N instead of 0.85\,N; at the wrist the motor inertia is about 100 times the link inertia.
  \item Turning the fork moves the tool tip: the law treats position and orientation as independent springs, but
        the arm's dynamics couple them. With the criterion (tip drift $\le 2$\,mm, angles $\le 2^\circ$) the first
        design limited rotations to 7.8$^\circ$/s. The acceleration limit (Section~\ref{sec:al}) now allows
        45$^\circ$/s within the criterion. Decoupling rotation and translation by inertia feedforward
        (operational-space control~\cite{khatib1987}) is option B3 of Section~\ref{sec:al}: planned only after a
        policy trains successfully and the robot's model has been checked (identification), because it makes the
        motion depend on the model.
  \item Close to a joint limit the arm lags more; there the real robot would stop. The task will keep the tool away
        from these poses (workspace bound or penalty).
  \item When pushing the stem, contact force (0.8 to 1.4\,N) and curvature (2.0 to 3.2\,1/m, limit 5) depend on
        the push, hardly on the controller: the stem is far softer than the controller's spring.
\end{itemize}

\paragraph{Limits.} Simulation only; robot parameters from Franka's description files, not identified on our
robot; no joint friction; control period 2\,ms instead of 1\,ms.

\paragraph{Questions.}
\begin{enumerate}
  \item Is this controller acceptable next to a plant ($K_p = 1000$\,N/m, $K_o = 100$\,N\,m/rad, at most 10\,cm/s and
        45$^\circ$/s, contact force bounded to 4\,N)?
  \item Lab setup: is the Franka Control Interface (torque control) enabled on our FR3, with which system and
        libfranka versions; is there a real-time PC for it; does the lab already use a control stack
        (franka\_ros2, frankapy, own C++) that this controller should be added to?
\end{enumerate}

{\raggedright  % long code names in the references: no stretched lines
\begin{thebibliography}{99}
\bibitem{tang2023industreal} B.~Tang, M.~A.~Lin, I.~Akinola, A.~Handa, G.~S.~Sukhatme, F.~Ramos, D.~Fox,
  Y.~Narang. IndustReal: transferring contact-rich assembly tasks from simulation to reality. \emph{RSS}, 2023.
\bibitem{frankatwin} FrankaTwin: aligned simulation and task-space impedance control for Franka robots in Isaac
  Lab. \texttt{github.com/tsrobcvai/frankatwin}.
\bibitem{kaspar2020sim2real} M.~Kaspar, J.~D.~Muñoz Osorio, J.~Bock. Sim2Real transfer for reinforcement learning
  without dynamics randomization. \emph{IEEE/RSJ IROS}, 2020.
\bibitem{peng2018sim} X.~B.~Peng, M.~Andrychowicz, W.~Zaremba, P.~Abbeel. Sim-to-real transfer of robotic control
  with dynamics randomization. \emph{IEEE ICRA}, 2018.
\bibitem{aljalbout2024action} E.~Aljalbout, F.~Frank, M.~Karl, P.~van der Smagt. On the role of the action space
  in robot manipulation learning and sim-to-real transfer. \emph{IEEE Robotics and Automation Letters}, 2024.
\bibitem{aljalbout2026realitygap} E.~Aljalbout, J.~Xing, A.~Romero, I.~Akinola, C.~R.~Garrett, E.~Heiden,
  A.~Gupta, T.~Hermans, Y.~Narang, D.~Fox, D.~Scaramuzza, F.~Ramos. The reality gap in robotics: challenges,
  solutions, and best practices. \emph{Annual Review of Control, Robotics, and Autonomous Systems}, 9, 2026.
\bibitem{albu2003cartesian} A.~Albu-Schäffer, C.~Ott, U.~Frese, G.~Hirzinger. Cartesian impedance control of
  redundant robots: recent results with the DLR-light-weight-arms. \emph{IEEE ICRA}, 2003.
\bibitem{khatib1987} O.~Khatib. A unified approach for motion and force control of robot manipulators: the
  operational space formulation. \emph{IEEE Journal on Robotics and Automation}, 3(1), 1987.
\bibitem{libfranka} Franka Robotics. libfranka example \texttt{cartesian\_impedance\_control.cpp}.
  \texttt{github.com/frankaemika/libfranka}.
\bibitem{frankaros} Franka Robotics. franka\_ros \texttt{cartesian\_impedance\_example\_controller}.
  \texttt{github.com/frankaemika/franka\_ros}.
\end{thebibliography}
}
% ======================================== END CONTENT ========================================
